\section{Further research questions and concrete next targets}
\label{sec:questions}

The proved results make the next questions more specific. The following
agenda concerns exact identities, structural reductions, and proof
certificates; it does not depend on searching for sharper inequalities.
Questions are labeled as questions, rather than as unproved theorems.

\subsection{Additional identities within the Dougall family}

\paragraph{1. Closed polynomial formulas for all resonant counterterms.}
The recurrence \eqref{dg:recurrence} is already finite and exact, and
\eqref{dg:quarticcoefficienttransport} builds every quartic moment
from one base sequence. Can its coefficients be written in a useful
closed combinatorial form in $N,j$, perhaps through central factorial
numbers and Bernoulli polynomials? A useful outcome would provide both
the subtraction polynomial and a compact exact certificate, avoiding
large expanded rational polynomials. This is a representation question;
the existence and the resonance value are proved here.

\paragraph{2. A minimal arithmetic presentation for rational-parameter jets.}
At rational $a,b,c$, Theorem~\ref{dg:alljets} gives a finite list of
Hurwitz derivatives and Stieltjes constants. Finite Fourier transforms
then express rational-shift Hurwitz values in residue-class or
Dirichlet-character coordinates. Which symmetries of a particular
Dougall family eliminate entire character sectors? The first concrete
case is the quartic family at shift $1/4$, comparing the Stieltjes
constants $\gamma_m(1/4)$ with the zero-order spectral jets
$\zeta^{(m)}(0,1/4)$ from its half-resonance. An exact reduction matrix
for each prescribed finite jet order would be useful. Numerical
independence is a separate question.

\paragraph{3. Mixed parameter jets along residue-zero loci.}
Formula~\eqref{dg:EN} makes all zeros explicit, but the article mainly
illustrates derivatives in $u$. When $b$ or $c$ is a positive integer
at most $N$, what short identities arise from mixed derivatives
transverse and tangent to $\rho_N=0$? The first target is a
two-parameter table at $b=c=1$, where the zero has higher multiplicity.
The finite product predicts which derivative orders first survive.
A complete statement should retain the small-index summands before
specialization, just as in \eqref{dg:revivalpi}.

\paragraph{4. Primitive endpoint values beyond the counterterms.}
The weighted primitive in Theorem~\ref{dp:primitive} lowers the
hypergeometric order and gives a finite colored-polylogarithm
description of the subtractions. It does not evaluate every endpoint
of the remaining \({}_4F_3\). Which parameter choices admit an
additional summation, a modular parametrization, or a finite
iterated-integral reduction? The squared-binomial branch is a
natural first case because its base generating function has
classical integral representations. Any proposed polylogarithmic
reduction should specify its function class and branches, rather
than infer it from the appearance of polylogarithmic counterterms.

\paragraph{5. Other exactly summable well-poised families.}
The present proof uses two special structures: reflected Gamma
shifts giving even powers, and an exact Gamma-product endpoint.
Which further very-well-poised summations retain both structures
under a one-parameter excess deformation? A precise first task is
to identify a higher-order family, derive its centered finite
asymptotic coefficients, and determine whether a differential
transport identity simplifies every resonant constant as in
\eqref{dg:closedeq}. The existence of the Dougall result does not
establish such a theorem for arbitrary hypergeometric families.

\subsection{Coordinate algebra and collisions}

\paragraph{6. Collision limits versus specified coincident finite parts.}
Let two translated singular points approach each other.
Corollary~\ref{nc:coincident} evaluates the coordinate change of the
finite part of the pointwise product at coincidence, while
Theorem~\ref{tr:finite-transport} treats disjoint supports.
What explicit local terms must be subtracted from the separated
family to obtain that coincident finite part? The first nontrivial
target is $\gamma_0(x)\gamma_0(x+\varepsilon)$, followed by
$\gamma_m^{(r)}(x)\gamma_n^{(s)}(x+\varepsilon)$.
The regulator and the relative approach direction must be fixed
before taking the limit. Equality of two unspecified
regularizations would not be a meaningful claim.

\paragraph{7. An intrinsic classification of the coordinate corrections.}
The residue law \eqref{nc:cocycle} is explicit. Can all local
corrections compatible with scalar pullback, finite principal parts,
and composition be classified, with the freedom of adding fixed
delta derivatives made precise? The sharp filtration
\eqref{nc:order-filtration} provides a concrete family of graded
invariants for that classification. A first step is to classify
the finite jet action through pole order three and determine
which coefficients can be changed by a fixed normalization.

\paragraph{8. Smooth coordinates and noninteger logarithmic exponents.}
For fixed pole order, the proof only uses finite coordinate jets
and an integrable remainder. This suggests an extension from
analytic coordinates to a stated finite differentiability class.
The exact minimum regularity and remainder hypotheses should be
proved, not inferred from formal Taylor manipulations. A distinct
question is the extension of $\log^n x$ to noninteger powers of
$\log x$, where branch choices and the polynomial filtration change.
It should be treated separately from the integer Stieltjes hierarchy.

\subsection{Colored kernels and exact relation certificates}

\paragraph{9. Minimal finite reductions for three or more frequencies.}
The finite lift has at most $p_1p_2p_3$ translated triples, and each
triple has six sign regions. Root filtering establishes existence
of a finite ordinary-kernel presentation. Can common divisors,
translation symmetry, and color conjugation compress this
presentation canonically? The first target is an exact reduction
algorithm for the $(1,2,3)$ example, with the order-derivative
coordinates retained. For four frequencies, one should specify
the resulting triple-sum kernels and a finite sign-region
decomposition before claiming analogous closure.

\paragraph{10. Transverse derivatives at negative Tornheim order.}
At $C=-N$, the kernel reduces to products of ordinary
polylogarithms. Equation~\eqref{tc:normal-derivative} isolates the
extra information in the first normal derivative. At roots of
unity, which choices of $A,B,N$ reduce that integral to a finite
combination of known depth-two polylogarithms, order derivatives,
and Gamma or Stieltjes data? The case $C=0$ with $A,B$ equal to
zero or one is a concrete starting point. The ordinary
negative-order value alone contains insufficient information.

\paragraph{11. Exact Gaussian $S_6$ and $S_8$ certificates.}
The retained conjectures need an exact relation in a specified
functional or iterated-integral system, together with endpoint
control and evaluation. The appropriate first task is to pin one
candidate vector, its alphabet, and the exact relation generators.
Either produce a certificate whose evaluation is the desired
identity, or prove exclusion from that \emph{specified} generated
relation space. Exclusion from a selected formal space would
identify a missing mechanism; it would not prove numerical
independence or disprove the period identity. The current source
ledger already emphasizes this distinction \cite{repo}.

\paragraph{12. A common exact coefficient and certificate format.}
The Bernoulli recurrence, the nonlinear residue, and the finite root
filter are all constructive operations on finite data once the
analytic theorem has been fixed. A useful ProveIt implementation
would record the theorem identifier, parameter domain, local
coordinate, coefficient order, finite-head terms, and exact
polynomial certificate in one reusable format. A first deliverable
could replay \eqref{dg:quarticN1jet},
\eqref{nc:r2unit}, and \eqref{tc:lower-polynomial} using
independent exact arithmetic. Such a format would support
integration and auditing without treating a floating-point
comparison as a proof certificate.

\subsection{Research priorities}
The shortest paths to additional exact identities appear to be
mixed Dougall jets at explicit residue zeros, arithmetic reduction
of the rational-shift special values, and normal derivatives of
colored Tornheim kernels at $C=0$. Collision limits require a
separate normalization theorem. The Gaussian conjectures remain
the most direct targets for a finite exact relation certificate,
but the present continuation supplies no reason to mark them
proved before that certificate exists.
